We describe \nonparametrics---\BSpline, \ac{BGP}, \ac{AR}, \vamana---used in this work. Results in the main text make use of \BSpline and \ac{BGP} models. \ac{BGP} was chosen since it was the only model capable of modeling the full mass spectrum, while \BSpline was chosen for most \ac{BBH} analyses for its ease of description and flexibility, as it simultaneously models all parameters with B-Splines. See Appendix~\ref{appendix:non_parameteric_comparison} for comparison of these \nonparametrics.
\subsection{\BSpline models} \label{appendix:BSpline}

\BSpline: The \BSpline model~\citep{Edelman:2022ydv} simultaneously fits all population parameters with Basis-splines (B-splines). 
A $k^\mathrm{th}$ order B-spline of variable $x$ consists of a linear combination of $n$ basis functions $\{B_{k,n}(x)\}$, each of which are degree $k-1$ piecewise polynomials joined at a set of $m$ locations called knots $\{x_m\}$. 
The number of basis functions $n$ defined across any range of $x$ values is determined solely by the order of the spline $k$ and the total number of knots $m$ as $n = k + m$.
This forms a basis that spans the space of possible interpolants between the knots $\{x_m\}$. Given a vector of coefficients $\pmb{\alpha}$, any function $f$ can then be approximated as 
\begin{equation}
\tilde{f}(x) = \sum^n_{i=1} B_{k,i}(x) \alpha_i.
\label{eqn:basis_splines}
\end{equation}
In the case of this work, $f$ is a probability distribution $p(\theta|\pmb{\alpha}_\theta)$ for a population-level parameter $\theta$, where the B-spline coefficients $\pmb{\alpha}_\theta$ are the hyperparameters of the distribution that are inferred during parameter estimation. 

Given an adequate number of knots, a B-spline is highly flexible and therefore capable of identifying sharp features that could be present in the data. 
This does, however, naturally make the \BSpline model prone to overfitting. 
To combat this, we include a smoothing prior $\pi_\text{BS}$ that penalizes large differences between neighboring coefficients. 
The prior is defined as 
\begin{equation}
p(\pmb{\alpha}|\tau) = \text{exp}\left( -\frac{1}{2} \tau \pmb{\alpha}^{\rm T} \pmb{D}_r^{\rm T} \pmb{D}_r\pmb{\alpha}\right),
\end{equation}
where $D_r$ is the $r$-order difference matrix with shape $(n - r) \times n$ and $\tau$ is a scalar that controls the level of smoothing. 
Ideally, $\tau$ is also inferred during parameter estimation. We found that $\tau$ consistently railed against prior boundaries, which meant that the limits imposed on the likelihood uncertainty was the main driver of smoothness. We therefore fix $\tau$ to a reasonable value (roughly between $5{-}10$) for each population distribution. 
With a sufficient number of bases, typically $n \sim 30{-}40$, this penalty prior will prevent the B-spline from overfitting the data while the large number of bases will provide enough flexibility to fit sharp features. 

The mass and spin distributions are modeled as B-splines, and the redshift distribution is modeled as a power law modulated by a B-spline \citep{Edelman:2022ydv}, with all components inferred simultaneously.
One of the \nonparametrics used in previous analyses~\citep{KAGRA:2021duu} was the \textsc{Powerlaw + Spline} model. This model assumed that the primary BBH mass followed a power-law distribution with moderate deviations controlled by a cubic spline.
The \BSpline model does not assume an underlying shape for the primary mass distribution, instead allowing for full model flexibility.

The \BSpline model infers the separable components of the mass distributions, $p(m_1)$ and $p(q)$, wherein the primary mass is defined over the range $3{-}300 \, \Msun$ and the mass ratio is defined over the range $0.03{-}1$. Unlike the \BrokenPLTwoPeaks model, a minimum secondary mass is not enforced during parameter estimation. To provide a more direct comparison to the \parametric, the \BSpline mass distributions shown in Section \ref{subsec:bbh_masses} are not the separable components $p(m_1)$ or $p(q)$ but instead the marginal distributions conditioned on $m_2 > 3 \, \Msun$, that is, $p(q|m_2 > 3 \, \Msun) = \int p(q)p(m_1)\Theta(m_1q-3)\mathrm{d}m_1$. We include the separable distributions along with the conditional marginal distributions in the mass ratio plot in Figure \ref{fig:bbh_mass_non_parametrics} to illustrate how this assumption affects the shape of the mass ratio distribution. 

\textsc{Isolated Peak}: This model is defined by two subpopulations~\citep{Godfrey:2023oxb}. One assumes a primary mass distribution described by a log-Gaussian peak while the other infers the mass distribution with a B-spline. Mass ratio, spin magnitude, and spin tilt distributions are inferred separately for each subpopulation, also using B-splines. The redshift distribution is the same for each subpopulation and is inferred with a power law modulated by a B-spline. 

\subsection{\acf{BGP} Model}\label{appendix:BGPModel}

In previous population analyses~\citep{KAGRA:2021duu}, we used the \ac{BGP} model to study the joint distribution of primary and secondary masses. Here, we extend it to model the joint distributions of mass, spin, and redshift.
The \ac{BGP} approach models the rate of \ac{BBH}, \ac{BNS} and \ac{NSBH} mergers as a piecewise constant function over a set of fixed bins across the one-, two-, or three-dimensional joint space. 
The comoving merger rate density in each bin is a hyperparameter of the model.
In addition, the \ac{BGP} model couples the logarithm of the rate density in each bin with a Gaussian process covariance, assuming an exponential quadratic kernel (also known as a radial basis function (RBF) kernel). 
The exponential quadratic kernel has a hyper-hyperparameter length scale $\lambda$ for each parameter in the joint space.

The covariance between two bins is proportional to the exponential of the negative squared distance between the bin centers, in units of the length scales along each parameter.
In addition, there is one more hyper-hyperparameter $\sigma$, which acts as an overall multiplicative scaling of the covariance matrix~\citep{Ray:2023upk}.
The \ac{BGP} approach directly infers the rate density in each bin, as well as the hyper-hyperparameters of the covariance kernel.
The prior on the parameter $\sigma$ is a half-normal with width $1$, and the length scale priors were tuned to the mean and variation in the set of distances between the bin centers.

We used 22 bins spaced uniformly in log-mass for the mass \ac{BGP} models, and 15 bins spaced uniformly between $\chi_{\rm eff} \in [-0.7,0.7]$ for the effective spin \ac{BGP} models.

The \ac{BGP} approach has the advantage of being able to model nontrivial correlations in the population of compact binaries. 
However, it assumes an arbitrary binning scheme, which reduces the resolution of the constraints and leads to unphysical discontinuities at the boundary between bins.
Furthermore, measurements of the rate density in some bins may be driven by the \textit{a priori} assumption of a Gaussian covariance, and not the data. This can also cause smoothing near sharp features e.g., near the $m_1=m_2$ boundary.

\subsection{\acf{AR} Model}

The \ac{AR} model \citep{Callister:2023tgi} is a highly flexible model for one dimensional marginal merger rate densities. 
The Monte Carlo integrals for estimating the likelihood in Equation~\eqref{eq:hierarchical_likelihood} involve a large set of parameter estimation samples [Equation~\eqref{eq:single_event_estimators}] and found injections for the selection efficiency [Equation~\eqref{eq:selection_estimator}].
The merger rate density at each sample is its own hyperparameter, which is directly inferred from the data. 
Without any further assumptions, such a model is severely underconstrained and will converge on the maximum likelihood functional distribution~\citep{Payne:2022xan}.
In order to \textit{a priori} favor smoother distributions, this finite list of samples is ordered along the dimension of interest like, for example, their primary mass. 
The marginal merger rate density is then assumed to be an autoregressive Gaussian random walk in log-space.
There are two hyper-hyperparameters, $\sigma_{\rm AR}$ and $\tau_{\rm AR}$, which control the scale of variability and the autocorrelation length along the Gaussian random walk respectively.
These hyper-hyperparameters have half-normal and log-normal priors respectively, tuned to the scale of the data: see Appendix B in \cite{Callister:2023tgi} for details. These hyper-hyperparameters are jointly inferred along with the logarithmic merger rate density at each sample.

The \ac{AR} model is particularly well-suited to distributions with sharp features and nontrivial evolution, complementing the \ac{BGP} approach. 
However, the \ac{AR} model used here cannot model correlations in the population. Furthermore, as in the \ac{BGP} approach, the inferred merger rate may extrapolate based off nearby constraints in regions of little information.
When evaluating uncertainties in the \ac{AR} model, one should consider the impact of the \ac{AR} smoothing prior. 
In particular, lower bounds on the merger rate in some regions may not be an accurate lower bound on the true distribution.

Like the \BSpline model, the \ac{AR} model does not enforce a minimum secondary mass through a mass-dependent mass ratio.

\subsection{\acf{FM} Model}

The \vamana model \citep{Tiwari:2020vym} is a flexible mixture model framework for modeling the \ac{BBH} population. 
\vamana models the population as a Gaussian mixture model in chirp mass and aligned spins, and a power law in mass ratio and redshift \citep{Tiwari:2021yvr}. 
Furthermore, \vamana is able to model correlations between parameters in the population, and additionally has a variable model dimension, sampling the trans-dimensional posterior using a reversible-jump MCMC method. For the analyses presented here, \vamana uses $11$ mixture components.
